\section{Verwacht resultaat, conclusie}%
\label{sec:verwachte_resultaten}

Naar verwachting levert dit onderzoek een algoritme op dat een e-mail analyseert, er een risicoscore aan toekent en de achterliggende redenen weergeeft via een lokale uitleg. Verwacht wordt dat het model op de testset een lage \emph{false-positive-rate} (FPR) behaalt. Tabel~\ref{tab:verwachte_resultaten} geeft een inschatting van de  prestaties weer.

\begin{table}[htbp]
  \centering
  \begin{tabular}{lc}
    \toprule
    \textbf{Metriek} & \textbf{Verwachting} \\
    \midrule
    FPR       & $\leqslant 3\%$ \\
    Recall    & $\geqslant 80\%$ \\
    Precision & $\geqslant 95\%$ \\
    F1-score  & $\approx 87\%$ \\
    Accuracy  & $\approx 88,5\%$ \\
    \bottomrule
  \end{tabular}
  \caption{Verwachte prestaties van het model op de testset.}
  \label{tab:verwachte_resultaten}
\end{table}

De FPR is de belangrijkste metriek (sectie~\ref{sec:evaluatie}) en moet dus zo laag mogelijk zijn. Het model wordt daarom bewust terughoudend afgesteld, waardoor een lagere \emph{recall} wordt aanvaard.

De \emph{precision} hangt af van het aandeel phishingmails in de testset. De verwachte waarde gaat uit van een evenwichtige verdeling van 50\% phishingmails en 50\% legitieme e-mails.

Verder is de F1-score het harmonisch gemiddelde van \emph{precision} en \emph{recall} en wordt de \emph{accuracy} ingeschat op ongeveer 88,5\%.

Wat betreft de kenmerken, wordt verwacht dat van de kandidaat-kenmerken uit sectie~\ref{sec:kenmerkextractie} een groot deel niet in het uiteindelijke model terechtkomt. Door de L1-regularisatie (sectie~\ref{subsubsec:machinelearning-detectie}) krijgen kenmerken met een beperkte bijdrage een gewicht van nul en vallen ze dus weg. Vermoedelijk wegen de kenmerken rond afzenderidentiteit en lexicale URL-kenmerken het zwaarst door.

Tot slot wordt verwacht dat elke bijdrage aan de risicoscore herleidbaar is tot controleerbaar bewijs in de e-mail. Voor elke e-mail met een risicoscore moet kunnen worden aangetoond welke onderdelen van het bericht hebben bijgedragen tot de eindscore. Zo wordt aangetoond dat een uitlegbare risicoscore technisch haalbaar is. Of zo'n uitleg gebruikers ook effectief helpt om phishingmails te herkennen, en of ze \emph{overreliance} vermijdt \autocite{Bucinca2021}, vraagt aanvullend gebruikersonderzoek en valt buiten de scope van dit onderzoek.